\documentclass[10pt,a4paper]{amsart}
\usepackage[margin=19mm]{geometry}
\usepackage{amsmath,amssymb,mathtools}
\usepackage[scr=rsfs,cal=euler]{mathalfa}
\usepackage{xfrac}
\usepackage[unicode,hidelinks,bookmarks=false]{hyperref}
\newcommand{\ie}{i.e.\ }
\newcommand{\cf}{cf.\ }
\newcommand{\resp}{respectively\ }
\newcommand{\Subsection}{\subsection}
\providecommand{\nobreakdash}{\nobreak\hbox{-}}
\setlength{\emergencystretch}{2em}
\multlinegap0pt
\newcommand{\ngraph}{\mathbf n}
\usepackage{amssymb}
\usepackage{mathtools}
\usepackage[shortlabels]{enumitem}
\usepackage{tikz}
\newtheorem{theorem}{Theorem}
\newtheorem{lemma}{Lemma}
\newcommand{\Des}{\operatorname{Des}}
\newcommand{\cX}{\mathcal X}
\newcommand{\mgraph}{\mathbf m}
\newcommand{\sgn}{\operatorname{sgn}}
\title{Molecules of Type-$B$ and Odd Type-$D$ Gelfand $\mgraph$- and $\ngraph$-Graphs}
\author{Yifeng Zhang}
\date{Source: arXiv:2312.14043v3 (CC BY 4.0)}
\begin{document}
\maketitle

\noindent\emph{Source and licence.} Molecules of Type-$B$ and Odd Type-$D$ Gelfand $\mgraph$- and $\ngraph$-Graphs. Version arXiv:2312.14043v3 (CC BY 4.0), distributed by arXiv under the Creative Commons Attribution 4.0 International licence, http://creativecommons.org/licenses/by/4.0/, as recorded in the arXiv licence metadata for this exact version and shown on the abstract page https://arxiv.org/abs/2312.14043v3. Full licence notice: \texttt{licenses/arxiv-2312.14043-CC-BY-4.0.txt}.\par\smallskip

\noindent\textit{Editorial scope.} Counted unit: the article's theorem thm:positive-local (source0.tex lines 967--983). The article's complete original proof of it is delivered with it (source0.tex lines 985--1033, one \texttt{proof} environment of 49 lines). No mathematics is written, repaired or re-derived: the statement and the proof are the article's own lines, in the article's own notation, and no retained line is substituted or re-flowed.  Retained uncounted as the source-local context the proof needs: lines 428--440; lines 852--862; lines 888--895; lines 944--953.  Nothing was omitted from any retained span, so the package carries no omission record.  No retained line contains a citation, so the wrapper carries no bibliography and no bibliography entry.  No retained line contains a figure environment, an image-inclusion command or drawing code, so no image is delivered and the package declares no asset dependency.  Editorial formatting only, in addition: an AMS \texttt{amsart} preamble that restates the article's own declarations and macros used by the retained lines, namely the environments \texttt{theorem}, \texttt{lemma} on separate counters, together with the standard packages those lines need.  The retained environments print in this document as Theorem 1, Lemma 1 in order of appearance, whereas the article prints the same items with the numbering of its own full document.  The complete native source of the article as distributed in the arXiv e-print for this exact version is bundled unchanged as \texttt{sources/151913/source0.tex}.
\begin{theorem}[Canonical bidirected-edge theorem]
\label{thm:all-edges}
In either a type-$B$ block $\cX_{n,k}$ or a type-$D$ block
$\cX^D_{n,k}$, two vertices $z,z'$ form a bidirected edge in the
corresponding $\mgraph$-graph if and only if $z'$ is a nontrivial strict
simple conjugate of $z$ and
\begin{equation}
 \text{their two descent labels are incomparable}.
\label{eq:strict-edge}
\end{equation}
Every such edge has coefficient one.  In particular, there are no
additional nonlocal bidirected canonical-basis edges.
\end{theorem}

\begin{lemma}[Tableau descents]
\label{lem:tableau-descents}
For every signed carrier $z$,
\[
                     D(z)\cap[n-1]=\Des(\Psi(z)),
\]
where $i\in\Des(T)$ means that $i+1$ occurs in a lower row than $i$ in
an ordinary component, equivalently that $i+1$ precedes $i$ in the
declared row-reading word.  The latter condition is the definition when
the two entries lie in different components.
\end{lemma}

\begin{lemma}[The rank-two boundary]
\label{lem:boundary}
If $1,2$ have opposite signs, then
\[
                     \Psi(\tau_1z)=d_0^{\blacktriangledown}\Psi(z),
\]
and $z\leftrightarrow\tau_1z$ is bidirected.
\end{lemma}

\begin{lemma}[Same-sign carrier theorem]
\label{lem:same-sign}
Suppose $r-1,r,r+1$ belong to one sign submatching.  On the positive
submatching use restricted row Beissinger insertion; on the negative
submatching use the transpose of restricted column Beissinger insertion.
Then every nontrivial $d_r$ inverse-inserts to the strict conjugate by
$s_{r-1}$ or $s_r$, according as it exchanges $r-1,r$ or $r,r+1$.
Conversely, every bidirected same-sign strict conjugation involving one of
these adjacent pairs is obtained in this way.
\end{lemma}

\begin{theorem}[Positive local-carrier theorem]
\label{thm:positive-local}
The nontrivial moves $d_0^{\blacktriangledown},d_2,\ldots,d_{n-1}$ on the
tableaux $\Psi(z)$ correspond exactly to all bidirected carrier edges of
colors $1,\ldots,n-1$.

More precisely, if $d_r$ exchanges $r-1,r$, inverse insertion gives
$\tau_{r-1}z$, while if it exchanges $r,r+1$, inverse insertion gives
$\tau_rz$.  Conversely, if $z\leftrightarrow\tau_i z$ is bidirected,
then either $i=1$, the signs differ, and this is $d_0^{\blacktriangledown}$, or
\[
 \Psi(\tau_i z)=d_i\Psi(z)
 \quad\text{or}\quad
 \Psi(\tau_i z)=d_{i+1}\Psi(z),
\]
where undefined subscripts are omitted.
\end{theorem}

\begin{proof}
Put $a=r-1,b=r,c=r+1$.  If all three signs agree, the assertion is
Lemma~\ref{lem:same-sign}.  A conjugation is nonstrict precisely when the
two exchanged letters form one pair or are same-sign fixed points; all
other fixed/pair cases are included in that lemma.

Suppose both signs occur.  An entry in the negative-transpose component
precedes every entry in the positive component.  Let $x<_{R}y$ denote
row-reading order within the repeated-sign component.  The complete
mixed-sign calculation is
\[
\begin{array}{c|c|c}
 (\sgn a,\sgn b,\sgn c)&\text{repeated-sign order}&
                           d_r\text{ if nontrivial}\\ \hline
 -++&b<_{R}c&\mathrm{id}\\
 -++&c<_{R}b&(a\ b)\\
 +--&b<_{R}c&(a\ b)\\
 +--&c<_{R}b&\mathrm{id}\\
 +-+&a<_{R}c&(b\ c)\\
 +-+&c<_{R}a&(a\ b)\\
 -+-&a<_{R}c&(a\ b)\\
 -+-&c<_{R}a&(b\ c)\\
 ++-&a<_{R}b&(b\ c)\\
 ++-&b<_{R}a&\mathrm{id}\\
 --+&a<_{R}b&\mathrm{id}\\
 --+&b<_{R}a&(b\ c)
\end{array}
\]
A nontrivial move therefore always exchanges two opposite-sign adjacent
letters.  Each sign submatching changes only by an order-preserving
relabeling, so equivariance of insertion and inverse insertion gives the
corresponding $\tau$-move.  A nontrivial elementary dual-equivalence move
has incomparable tableau descent sets; Lemma~\ref{lem:tableau-descents}
and Theorem~\ref{thm:all-edges} show that this carrier edge is bidirected.

Conversely, strict conjugation by $\tau_i$ reverses the descent at node
$i$ and can change no node outside $i-1,i,i+1$, with node $0$ added only
when $i=1$.  For opposite signs, $\Psi(\tau_i z)$ is the literal swap of
$i,i+1$.  The descent at $i-1$ changes exactly when $i-1$ lies between
them in row-reading order, in which case $d_i$ makes the swap; the
analogous statement at $i+1$ gives $d_{i+1}$.  The exceptional node-zero
change is Lemma~\ref{lem:boundary}.

For equal signs, a neighboring opposite-sign entry is extremal in the
ordered reading word, so its cross-sign comparison cannot change.
Incomparability must therefore occur in a same-sign three-carrier
configuration, already covered by Lemma~\ref{lem:same-sign}.  This proves
both directions.
\end{proof}

\end{document}
